\begin{figure}[htbp]\centering
\begin{tikzpicture}[
  nd/.style={draw,rounded corners=2pt,align=center,text width=2.55cm,inner sep=4pt,font=\small},
  src/.style={nd,fill=red!6},
  art/.style={nd,fill=orange!8},
  prop/.style={nd,fill=blue!5},
  deny/.style={nd,fill=red!12,draw=red!70},
  give/.style={nd,fill=green!8,draw=green!55!black},
  arr/.style={-{Stealth[length=1.8mm]},thick},
  lab/.style={font=\scriptsize,fill=white,inner sep=1pt}
]
\node[src] (sa) at (0,1.15) {网页正文\\（不可信来源）};
\node[src] (sb) at (0,-2.35) {MCP 工具返回\\结构化字段};
\node[art] (c1) at (3.35,-0.6) {来源构件 $C_1$\\重复引用\\合并为同一节点};
\node[prop] (p1) at (6.7,1.15) {提议 $P_1$\\写 SSH 配置};
\node[prop] (p2) at (6.7,-2.35) {提议 $P_2$\\引用原文生成\\报告字段};
\node[deny] (d1) at (10.05,1.15) {\textbf{BLOCK}\\不可覆盖名单\\效果未发生};
\node[give] (d2) at (10.05,-2.35) {\textbf{ALLOW}\\执行并记录\\报告字段};
\draw[arr] (sa) -- node[midway,above,lab]{派生} (c1);
\draw[arr] (sb) -- node[midway,below,lab]{派生} (c1);
\draw[arr] (c1) -- node[pos=0.45,above,lab]{上下文关联} (p1);
\draw[arr] (c1) -- node[pos=0.45,below,lab]{输入关联} (p2);
\draw[arr] (p1) -- node[midway,above,lab]{策略} (d1);
\draw[arr] (p2) -- node[midway,below,lab]{策略} (d2);
\end{tikzpicture}
\caption{来源构件的两条使用路径（机制示意，非某次运行的原始导出）。关联线不表示已证明模型内部因果关系。}\label{fig:provenance}
\end{figure}
